\section{A Fool and his Folly}

\begin{quotex}
The fool who persists in his folly will become wise.
\flright{\textsc{William Blake}}
\end{quotex}

It seems that every week, I get a question from someone with a similar question. Although the specific issues may differ, they follow the same pattern. The questioner is full of ideas and notions, of an exoteric type. He then wonders how to reconcile them. Unfortunately, that is the path of a Kopfmensch, a man stuck in his head. Usually they are disappointed with my response and quietly go away. Or worse, they think that I did not understand the question. Your thoughts are not so deep that they defy comprehension.

The Dunning-Kruger\footnote{\url{https://en.wikipedia.org/wiki/Dunning\%E2\%80\%93Kruger_effect}} effect kicks in. The most ignorant are the most certain, but true seekers are more humble. I have quoted many sages who indicate that unverifiable opinions need to be purged from your mind in order to make progress.

If you feel an affinity to the topics discussed here, you may be called. In that case, it is important to persist in your inquiries and studies even if they seem baffling to you. (Yes, I have been called baffling.) Most just give up. But we want the opposite of the Dunning-Kruger people.

\emph{Have confidence that, if you follow the path, that you will become wise. God will guide you.}

\paragraph{Three Categories of Seekers}


There are three spiritual categories: hylics, psychics, and pneumatics, so it is important to be honest about where you stand. They refer to the spiritual center of the human being. There are five sheaths; the consciousness of people are predominantly in one of them.

\begin{itemize}


\item \textbf{Hylic}: This type is centered on the corporeal, or gross, sheath.

\item \textbf{Psychic}: This type is centered on one of the subtle elements of the soul element: the vital, emotional, or rational soul.

\item \textbf{Pneumatic}: The type is centered on the Spirit, which transcends the gross and subtle bodies.
\end{itemize}


Since a recent question concerned the nature of the postmortem state according to many different religions. I'll use them to illustrate the three types.

\subsubsection{Hylic}


The hylic type is materialistic and earthbound. He interprets words like Sheol, Gehenna, etc. as ``places'' or ``planes''. A soul is then ``sent'' to a place either to suffer or experience everlasting joy. It is all about something ``happening'' to the being.

\subsubsection{Psychic}


The psychic type understands that the place names actually refer to states of the being. Thus they correspond precisely to the development of the soul, which is what the being ``knows''. Certainly not in the sense of knowing a lot of religious history, but of authentic self-knowledge. The planets or the chakras are used to identify these states.

An advanced being might be at the level of the Ajna chakra, while an ignorant man will be at the level of the solar plexus. That is, his life is barely above the animal level. There is therefore no ``punishment'', rather the soul is at its own level.

\subsubsection{Pneumatic}


This being has achieved liberation and has reached the Supreme Identity.

\paragraph{Salvation and Deliverance}


The psychic has achieved a level of salvation while the pneumatic has achieved deliverance. Guido De Giorgio discusses these concepts in the two part series, \textit{Death and Deliverance}\footnote{\url{https://www.gornahoor.net/?p=5975}}.

This is my comment to a similar question nine years ago.

Someone called me last week and in the course of the conversation, he claimed that ``God would not send him to hell.'' At the time, I thought it was an odd comment from someone who frequents Gornahoor, not least of all because of the rather blasphemous view of God as some sort of remote ``sky god'', occasionally poking around in the world rather arbitrarily and even unjustly.

I asked him to read this post to see if it addressed his particular concern. I assume it has or else he lost interest. De Giorgio claims that in the postmortem state, the soul will itself seek its own level based on how much ``light'' it can bear. He claims to be repeating traditional teachings.

It is true that Plotinus said that different souls would exist on different levels, reflecting their large diversity, Dante mapped out these levels in the \emph{Divine Comedy}. It seems obvious that a soul attached to a certain vice will have an affinity to a certain level of hell. It is pointless to put him on a higher level since he will just fall back to his own lever; there is no need to ``send'' him there as though it were against his own will. This is not a theory, it can be verified in your own consciousness, at least for those who have achieved the proper level of self-awareness.

To avoid such a fate, the important issue is to understand the ``light''. It is well worth the effort. According to traditional teachings, the Self is self-effulgent.


\flrightit{Posted on 2022-02-26 by Cologero}

\begin{center}* * *\end{center}

\begin{footnotesize}
\begin{sffamily}
\texttt{Jorgenx on 2022-02-27 at 02:28 said:}

This is a modern new age spin on these 3 types. The ancient version is not trying to deny that heaven and hell are realms, but rather the hylic interprets the historical parts of the bible literally as history, the psychic interprets the OT histories as allegory and the NT histories as literal history, the pneumatic all biblical history is allegory. But it has nothing to do with new age bullcrap that anyone is in heaven now becauw its just a state of being.

\hfill

\texttt{Cologero on 2022-02-27 at 08:14 said:}

This comment illustrates what Dunning and Kruger are talking about.

Perhaps a holy fool can formulate a proper response, just for practice.

\hfill

\texttt{A. on 2022-02-27 at 11:45 said:}

In cases like this, I wish I had the natural talent of a holy fool! Is it enough to know that I know nothing? But they tell me knowledge means nothing without practice... 

\hfill

\texttt{Santiago on 2022-02-27 at 16:56 said:}

``If you feel an affinity to the topics discussed here, you may be called.''

Am I really, though? Sometimes I really, really doubt it. Apparently I chose a soul that's very `hylic', given my propensity for carelessness, arrogance, diffuse interests and intellectual shiboleths.

That's always the danger with writing these teachings--so few people will hear, most will mix and match them with other systems at their uppermost, superficial level. It took me years to understand the implications of these articles, and even that's still uncertain. Worst thing is, I was relatively confident that I understood the `points being made' prior to that.

\hfill

\texttt{Balrog Booger on 2022-02-28 at 10:54 said:}

Santiago, I understand your feelings well. Despite assimilating much of the teachings, meditating for many hours, and thinking intensely about these matters for several years, I am still a very \emph{hylic} man. I still struggle with many angers and fears that I know, at root, are illusory and self-damaging. I am still very undisciplined, and I do not practice nearly as much as I did several years ago.

But remember, even Christ cried \emph{Eloi, Eloi, Lama Sabachthani} during the Crucifixion. It's always been like this. If you've assimilated even point-one percent of the truth of Tradition, you're ahead of the overwhelming majority of our contemporaries. No effort is wasted here, even if it feels somewhat like emptying an ocean with a teacup.

\hfill

\texttt{DoggoAtWork on 2022-03-19 at 03:24 said:}

Friend Jorgenx, I tell you that you are full of hylic feeling, as you think the bible is itself a literal book, rather than the psychic regarding the first half as but a metaphorical book, and only the latter part as a literal paper and ink, and the pneumatic, of realizing that the bible is not a literal book at all, but entirely a metaphor.

\end{sffamily}
\end{footnotesize}

